\subsection*{Impactos ambientales}
\addcontentsline{toc}{subsection}{Impactos ambientales}

El modelo operativo de la plataforma, al ser íntegramente digital, genera una serie de impactos ambientales positivos derivados de la desmaterialización de procesos y la optimización de recursos. La transición de la educación financiera tradicional hacia un entorno virtual elimina la dependencia de insumos físicos y materiales impresos, lo que reduce significativamente la generación de residuos sólidos y el consumo de papel. Asimismo, el despliegue de canales de soporte y comunicación síncrona mediante la plataforma contribuye a la reducción de la huella de carbono por movilidad, al eliminar la necesidad de desplazamientos físicos de los usuarios hacia centros de asesoría, disminuyendo así las emisiones de CO2 vinculadas al transporte urbano en Bogotá.

No obstante, se reconoce que el funcionamiento de la infraestructura cloud, los microservicios y el procesamiento masivo de datos generan una demanda energética que incide en la huella de carbono digital. Según \textcite{Belkhir2018}, se proyecta que las emisiones derivadas del sector de las Tecnologías de la Información y las Comunicaciones (TIC) representen una proporción creciente del inventario global de gases de efecto invernadero si no se implementan estrategias de mitigación. Para contrarrestar este efecto, el proyecto contempla la selección estratégica de proveedores de centros de datos que implementen políticas de energía renovable y estándares de eficiencia energética (PUE), alineándose con las prácticas globales de \textit{Green Computing}.

Finalmente, la plataforma genera un impacto indirecto pero estructural al promover comportamientos de consumo responsable entre sus usuarios. Al fomentar la planificación financiera frente al consumo impulsivo, se induce una modificación en los patrones de compra hacia decisiones más reflexivas y eficientes. Este cambio de mentalidad no solo fortalece la economía personal, sino que reduce la presión ambiental derivada de la producción masiva de bienes no esenciales, promoviendo una cultura de consumo alineada con la preservación de recursos naturales a largo plazo.